\section{保半正定的核装配与 BFZI 在线计算}\label{sec:positive-assembly}
\subsection{瓶颈和目标}
第 \ref{sec:admissible-atlas} 节的局部图谱减少了
跨几何重建 FEM 网格的需求，但 \(\bm R,\bm L,\bm K\) 的空间积分
仍可能成为推理瓶颈。因此需要在\textbf{保持物理 Gram 结构}的前提下
改造核装配。下面区分两种来源：
\textbf{精确对称积分 / 一致离散}和
\textbf{用正权核特征近似后仍具有半正定结构}。
不能通过独立预测矩阵元素来保证整体功率守恒。

\subsection{电阻 Gram：正权局部求积}
对同一套几何条件化实电流基 \(\bm W_i\)，
在实际导电域选取正权求积
\((\bm x_\nu,w_\nu)\)，\(w_\nu>0\)，并定义
\begin{equation}
\widehat{\bm R}_{ij}
=\sum_{\nu=1}^{Q_R}
\frac{w_\nu}{\sigma_s(\bm x_\nu,T)}
\bm W_i(\bm x_\nu)\cdot\bm W_j(\bm x_\nu).
\label{eq:R-positive}
\end{equation}
若用 \(\bm D_R\) 表示把每个模式在所有求积点的三分量
叠加后形成的加权采样矩阵，则
\[
\widehat{\bm R}=\bm D_R^{\mathsf T}\bm D_R\succeq0.
\]
若在采样点上字典列仍线性独立，则 \(\widehat{\bm R}\succ0\)；
若不独立，应消除冗余模态、细化近界面求积或选择稳定秩，
而不是在未解释空间损耗的情况下随意添加对角电阻。
用\textbf{同一求积节点}的
\(q_{s,\nu}=|\sum_i z_i\bm W_i(\bm x_\nu)|^2/(2\sigma_s)\)，
可保证
\[
\tfrac12\bm z^{\mathrm H}\widehat{\bm R}\bm z
=\sum_\nu w_\nu q_{s,\nu}.
\]
这只是\textbf{离散求积功率等式}。
如果用其它积分规则或连续精确热积分，
误差由 \(\widehat{\bm R}-\bm R\) 控制，
不得把离散等式当成零误差的连续物理结果。

\subsection{磁核 Gram：共同的非负 Fourier 特征}
在均匀自由空间磁导率、场足够衰减且基函数可作 Fourier 变换时，
记 \(\widehat{\bm W}_m(\bm k)=
\int_{\mathbb R^3}\bm W_m(\bm x)e^{-\jj\bm k\cdot\bm x}\,\mathrm d\bm x\)。
对 Coulomb Green 核
\(G_0(\bm x)=1/(4\pi\|\bm x\|)\)
有形式上
\[
\widehat{G_0}(\bm k)=\|\bm k\|^{-2}.
\]
在满足积分收敛的条件下，
由 Parseval 恒等式
\begin{equation}
[\bm L]_{mn}
=\frac{\mu_0}{(2\pi)^3}
\int_{\mathbb R^3}
\frac{\operatorname{Re}[
\widehat{\bm W}_m(\bm k)^*\cdot
\widehat{\bm W}_n(\bm k)]}
{\|\bm k\|^2}\,\mathrm d\bm k.
\label{eq:magnetic-fourier}
\end{equation}
该表达在 \(\bm k=0\) 附近的可积性、
高频收敛和外海水模式必须由实际函数空间保证。
用排除零波数奇点并匹配正权的 Fourier 节点
\((\bm k_\nu,a_\nu)\)、\(a_\nu>0\)，定义
\begin{equation}
\boxed{
\widehat{\bm L}_{mn}=
\mu_0\sum_{\nu=1}^{Q_L}
\frac{a_\nu}{\|\bm k_\nu\|^2}
\operatorname{Re}\!\left[
\widehat{\bm W}_m(\bm k_\nu)^*\cdot
\widehat{\bm W}_n(\bm k_\nu)\right].
}
\label{eq:L-positive}
\end{equation}
对于所有实系数向量 \(\bm c\)，
\[
\bm c^{\mathsf T}\widehat{\bm L}\bm c
=\mu_0\sum_\nu
\frac{a_\nu}{\|\bm k_\nu\|^2}
\left\|\sum_m c_m\widehat{\bm W}_m(\bm k_\nu)\right\|^2
\ge0.
\]
因此只要所有模式使用\textbf{同一组正权谱特征}，
即使求积尚有误差，所得实对称磁矩阵仍半正定。
实际实现可用带 \(\pm\bm k_\nu\) 成对的实/虚特征
保证矩阵实对称；\textbf{此构造是标准正定核性质的具体应用，
不是首次发现 Fourier 保正性的数学定理。}

对铜载流提升场、盒壁模式和海水模式，
可把全部电流函数列组成 \(\bm V=[\bm U,\bm W]\)，
用完全相同的 Fourier 特征生成完整
\(\widehat{\bm L}_{ab}\) 块；
这样\textbf{全部块联合半正定}，
联合 Schur 补的被动性和离散功率证明继续适用。
若海水 \(\bm K\) 使用精确 BFZI 线积分、
\(\bm L_{00}\) 使用另一组独立近似，
则不能自动断言合并后的完整磁 Gram 为半正定；
必须核对共同能量内积或回到联合同源组装。

\subsection{刚体变换可复用的谱特征}
对随组件刚体移动的模式
\(\bm W_p(\bm x)
=R_p\widehat{\bm W}_p(R_p^{\mathsf T}(\bm x-\bm c_p))\)，
其 Fourier 变换严格满足
\begin{equation}
\boxed{
\widehat{\bm W}_p(\bm k)
=e^{-\jj\bm k\cdot\bm c_p}\,
R_p\widehat{\widehat{\bm W}}_p(R_p^{\mathsf T}\bm k).
}
\label{eq:rigid-fourier}
\end{equation}
这里 \(\widehat{\widehat{\bm W}}\) 表示\textbf{局部参考模式的 Fourier 变换}，
不表示再对 Fourier 数据作一次变换。
因此局部几何固定且其模式不侵犯另一实体时，
只需改变相位、方向和频点采样，就能更新相对位姿的谱 Gram。
在任意旋转下，参考谱需要在旋转后的
\(R_p^{\mathsf T}\bm k_\nu\) 处评估；
插值误差可改变数值精度，
但若最终 \(\widehat{\bm L}\) 仍由同一组采样特征
形成 Gram，半正定性仍得到保持。
相对位姿变化不会使跨组件磁耦合不变；
\(\bm c_2-\bm c_1\) 对应的相位仍需计算。

\subsection{BFZI 快速构造端口耦合 \(\bm K\)}
第 \ref{sec:galerkin} 节的
\(\bm K_{mp}=\mu_0\int_{\Omega_s}\bm W_m\cdot\bm C_p\,dV\)
由 Green 核互易性可改写为
\begin{equation}
\boxed{
K_{mp}=\int_{\Gamma_p}\bm A_m^{\rm mode}(\bm r_p)
\cdot \mathrm d\bm\ell_p,\quad
\bm A_m^{\rm mode}(\bm r)
=\mu_0\int_{\Omega_s}G_0(\bm r-\bm x)\bm W_m(\bm x)\,\mathrm dV.
}
\label{eq:K-reciprocal}
\end{equation}
这表明可以\textbf{预计算或快速查询电流模式的矢量势}，
再沿 BFZI 连续线圈执行结构化线积分，
不必在每次端口查询中完整求一个海水三重体积分。
在固定线段组与 Gauss 节点，BFZI 仍满足
\(\bm r_{b,k}=\bm r_{b,0}+k\Delta\bm r_b\) 和
\(\bm h_{b,k}=\bm h_{b,0}+k\Delta\bm h_b\)。
若 \(\bm A_m^{\rm mode}\) 在匝号块覆盖的路径邻域
具有可控导数，可按 \(k\) 做 Taylor / Chebyshev 模式求和，
用
\(\sum_{k\in T}(k-k_c)^j\) 的显式矩求和加速；
在近场陡变、靠近盒边和模式奇异位置回退直接求积。
注意这与原 BFZI 论文对 \(1/\|\bm x-\bm r_k\|\) 核做二阶块求和
是\textbf{扩展应用}，新的 \(\bm A_m^{\rm mode}\) 误差界需单独推导，
不能沿用原文的余项常数。

\subsection{正定性、连续损耗与加速之间的取舍}
\textbf{路径 A（首选保结构）}：
所有 \(\bm R\) 均来自实际分区上的正权积分，
所有 \(\bm L\) 块来自同一正权谱/空间核 Gram；
因此任意子空间选择与联合 Schur 补保持离散被动性，
输出空间 \(q_i(\bm x)\) 由同一当前电流模式求得。
\textbf{路径 B（严格源核）}：
BFZI 互易线积分快速给出精确/有界的 \(\bm K\)。
此时对只有海水的低维散射阻抗，独立
\(\bm R_s\succ0,\bm L_s\succeq0\) 足以保证
海水散射部分被动；但若要把铜/盒壁共同耦合进完整 Schur，
则必须继续检查\textbf{完整磁 Gram 的一致性}。
\textbf{禁用的捷径}：
单独神经预测 PSD \(\bm R\) 和各点损耗却不给出相同的场内积，
不能宣称端口积分功率严格相等。

\subsection{在线复杂度的正确统计}
若字典模式数 \(M\)、学习子空间维数 \(r\)、电阻求积点数
\(Q_R\)、谱节点数 \(Q_L\)，且
\(\bm D_R,\bm D_L\) 可在当前场景生成，
则可先计算小型特征
\[
\bm F_{R,r}=\bm D_R\bm S_\theta,\qquad
\bm F_{L,r}=\bm D_L\bm S_\theta,
\quad
\bm R_r=\bm F_{R,r}^{\mathsf T}\bm F_{R,r},\quad
\bm L_r=\bm F_{L,r}^{\mathsf T}\bm F_{L,r},
\]
从而避免显式构造 \(M\times M\) 全矩阵。
典型稠密运算的贡献为
\(O((Q_R+Q_L)Mr+(Q_R+Q_L)r^2+r^3)\)，
\textbf{不包含}字典建立、特征值评估、
BFZI \(\bm K\) 核装配及最终场查询。
若分区紧支撑，\(\bm D_R\) 实际稀疏，
应利用局部结构减少常数和矩阵乘法代价。
要获得真实在线加速，必须使
\(Q_R,Q_L,M\) 与图谱计算在目标参数域内有可复用/分离结构；
仅令 \(r\) 很小并不足够。

\subsection{几何重组时的失败回退}
相对位姿导致局部模式支撑侵犯另一盒体、
间隙接近零、正权采样不足、低阶电流基不满秩或
谱/局部核逼近无法达到容差时：
增广间隙图谱、加密局部求积、缩小模型适用域
或回退高精度场求解。
不能为维持速度而输出不合规的电流或负功率损耗。
